\documentclass[12pt]{article}
\usepackage{amsmath}
\usepackage{amssymb}
\usepackage{hyperref}

\title{Evaluating Parameter Copying vs. Iterative Teaching in Matrix Environments}
\author{Prince T. Philip \\ \href{mailto:princephilip514@gmail.com}{princephilip514@gmail.com}}
\date{\today}

\begin{document}

\maketitle

\begin{abstract}
In this paper, we evaluate the computational efficiency of iterative cultural transmission (teaching) versus exact parameter copying in simulated intelligent systems. Utilizing a perfectly synchronized dual-RNG artificial life simulation (Matrix) to eliminate stochastic confounds, we compared a specific 3-step active error-correcting teaching loop (C-full) against direct parameter mimicry (B-copy). Our empirical results demonstrate that exact copying achieves a strictly superior task success rate over the specified iterative teaching method (88.42\% vs 84.73\%, $p_{\mathrm{MC}} < 10^{-5}$). However, this finding is largely bounded by the arbitrary hyperparameter limits of the teaching protocol. We conclude that exact copying trivially wins when transmission is treated as computationally free and frictionless. The true theoretical frontier for generational artificial life lies in exploring environments where copying incurs a penalty, parent/child architectures differ, or rapidly shifting environments render exact parameter clones obsolete.
\end{abstract}

\section{Introduction}
The emergence of intergenerational knowledge transfer (culture) is widely considered a hallmark of advanced intelligence. Biological teaching---a feedback loop of demonstration, attempt, and correction---is a highly complex behavior.

With the advent of artificial intelligence, we must ask: If a system is capable of instant, perfect parameter transmission (copying), does an algorithmic construct of ``teaching'' remain advantageous? This study mathematically models this dynamic for digital agents within a constrained framework. 

\section{Methodology}
To prevent global RNG contamination, we deployed a \textbf{Dual-RNG Architecture}. Each agent instantiated two separately seeded deterministic pseudorandom streams: an \texttt{inheritance\_rng} and a \texttt{lifetime\_rng}. This guaranteed that all populations experienced identical lifetime environmental challenges.

The inheritance mechanisms were defined as:
\begin{itemize}
    \item \textbf{Condition B-copy (Exact Copy):} Offspring receive a direct, uncorrected mathematical duplication of the parent's skill matrices and semantic concepts.
    \item \textbf{Condition C-full (Iterative Teaching):} Offspring undergo a simulated teaching loop. The algorithm calculates the error between the child's initialization and the parent's capability, applying exactly three iterative, stochastic corrections.
    \item \textbf{Condition A (No Inheritance):} Offspring start by initializing skills to zero and concepts to an empty set.
\end{itemize}

Populations of 100 agents were simulated over 10 generations. The environment was processed across 100 independent universe seeds.

\section{Results}
The empirical data shows a divergence in task success rates by Generation 10:
\begin{itemize}
    \item \textbf{B-copy (Exact Copy):} $\mu = 0.8842$ (88.42\%)
    \item \textbf{C-full (Iterative Teaching):} $\mu = 0.8473$ (84.73\%)
\end{itemize}

A paired sign-flip permutation test (100,000 permutations) on the 100 seeds yielded an empirical probability of $p_{\mathrm{MC}} \approx 10^{-5}$, indicating strong statistical divergence under these specific constraints.

\section{Discussion \& Limitations}
While B-copy numerically outperformed C-full, careful independent validation reveals that this outcome is heavily determined by construction. 

In C-full, ``teaching'' is explicitly bounded to three noisy steps toward the parent's skill value. As a result, C-full mathematically cannot overshoot the parent's value, and therefore exact copying (B-copy) will consistently tie or win. The measured performance gap is essentially an artifact of the 3-step hyperparameter. 

Furthermore, the simulation treats skills as capping at 1.0, meaning generations act functionally as a longer continuous lifespan rather than an interactive cultural mechanism. 

The interesting scientific challenge remains untested: What occurs when copying is not ``free''? In realistic physical or advanced neural architectures, copying carries substantial computational, storage, or energy costs. Furthermore, in rapidly non-stationary environments, a perfect clone of a parent's neural weights may inherently encode obsolete strategies, granting an advantage to iterative, generalized teaching methods that compress knowledge rather than memorizing state.

\section{Conclusion}
Under the tested simulation conditions, exact-copy inheritance achieved a higher mean final task-success score than the specified 3-step stochastic iterative error-correction inheritance method across 100 paired seeds. Future research must discard frictionless copying assumptions and introduce lifetime-matched single-agent baselines, architectural transmission costs, and shifting environmental parameters to truly probe the boundary between copying and teaching.

\end{document}
